\section*{Capítulo \ref{ch:g9:climate-past-and-present} --- El clima de la Tierra, ayer y hoy}

\begin{solution}{exo:g9:climate-past-and-present:1}
(a) Unos \qty{-19}{\degreeCelsius}. (b) Unos \qty{15}{\degreeCelsius}.
(c) $15 - (-19) = 34$: unos \qty{34}{\degreeCelsius}. (d) Dos de estos:
dióxido de carbono, vapor de agua, metano.
\end{solution}

\begin{solution}{exo:g9:climate-past-and-present:2}
(a) \omterm{def:g1:watching-the-weather:weather}{Tiempo}. (b) \omterm{def:g5:weather-and-climate:climate}{Clima}. (c) \omterm{def:g5:weather-and-climate:climate}{Clima} (y \omterm{def:g9:climate-past-and-present:climate-change}{cambio climático}: un cambio a lo largo
de décadas). (d) \omterm{def:g1:watching-the-weather:weather}{Tiempo}.
\end{solution}

\begin{solution}{exo:g9:climate-past-and-present:3}
(a) \qty{19}{cm}, once años. (b) $19 \div 11 \approx \qty{1.7}{cm}$.
(c) La \omterm{def:g4:clouds-rain-snow:precipitation}{nieve} de cada año cae encima de la \omterm{def:g4:clouds-rain-snow:precipitation}{nieve} de los años anteriores, así
que cuanto más profunda es una capa, antes cayó.
\end{solution}

\begin{solution}{exo:g9:climate-past-and-present:4}
(a) $22.4 - 21.0 = +\qty{1.4}{\degreeCelsius}$. (b) $20.3 - 21.0 =
\qty{-0.7}{\degreeCelsius}$. (c) \qty{0}{\degreeCelsius}. (d) Que ese mes
fue más frío que su media de 1991--2020.
\end{solution}

\begin{solution}{exo:g9:climate-past-and-present:5}
(a) Las décadas de 1880 a 1930, la de 1950 y la de 1960. (b) Unos
\qty{0.25}{\degreeCelsius} y \qty{0.80}{\degreeCelsius}. (c) $0.80 -
0.25 = \qty{0.55}{\degreeCelsius}$. (d) No: las décadas de 1880 a 1910 se
enfriaron un poco, y la de 1950 fue algo más fría que la de 1940; pero
desde la década de 1970 cada década ha sido más cálida que la anterior.
\end{solution}

\begin{solution}{exo:g9:climate-past-and-present:6}
(a) $(324.62 - 315.98) \div 10 \approx \qty{0.86}{ppm}$ al año. (b)
$(427.35 - 401.01) \div 10 \approx \qty{2.63}{ppm}$ al año. (c) $2.63
\div 0.86 \approx 3.1$: unas tres veces más rápida. (d) $(427.35 - 315.98)
\div 315.98 \times 100 \approx 35~\%$.
\end{solution}

\begin{solution}{exo:g9:climate-past-and-present:7}
(a) $8 \times 10^5$ años. (b) Hace unos \num{120000} a \num{130000} años:
la curva de temperatura sube por encima de 0 en esa época. (c) Unas
\qty{275}{ppm}. (d) Sí: los picos y los \omterm{def:g1:landscapes-around-us:landform}{valles} de las dos curvas aparecen
en los mismos momentos.
\end{solution}

\begin{solution}{exo:g9:climate-past-and-present:8}
(a) $3.2 \times 10 = \qty{32}{mm} = \qty{3.2}{cm}$. (b) $2100 - 2025 =
75$ años; $75 \times 3.2 = \qty{240}{mm}$, unos \qty{24}{cm} (y más si la
subida se acelera). (c) El agua del \omterm{def:g2:water-on-earth:ocean}{mar} se dilata al calentarse; el hielo
continental se funde y su agua va a parar al \omterm{def:g2:water-on-earth:ocean}{mar}. (d) No: el \omterm{def:g8:ice-on-earth:sea-ice}{hielo marino}
ya flota en el \omterm{def:g2:water-on-earth:ocean}{mar}, así que su fusión no lo hace subir.
\end{solution}

\begin{solution}{exo:g9:climate-past-and-present:9}
(a) 90. (b) 66. (c) 99. (d) \emph{Probable} significa al menos 66
posibilidades sobre 100, no certeza: el periódico presenta la afirmación
como más segura de lo que dijeron los científicos.
\end{solution}

\begin{solution}{exo:g9:climate-past-and-present:10}
Un invierno frío en un lugar es \emph{\omterm{def:g1:watching-the-weather:weather}{tiempo}}: dice lo que ocurrió durante
unos pocos meses en un punto. El \omterm{def:g9:climate-past-and-present:global-warming}{calentamiento global} es un cambio de
\emph{\omterm{def:g5:weather-and-climate:climate}{clima}}, medido a lo largo de décadas, y \emph{global}: es la media
de toda la Tierra. Un invierno frío aquí puede darse mientras la Tierra en
su conjunto está más caliente que antes, y las medias por décadas siguen
subiendo.
\end{solution}

\begin{solution}{exo:g9:climate-past-and-present:11}
(a) $280 - 190 = \qty{90}{ppm}$, en unos $2 \times 10^4$ años. (b) $427
- 280 = \qty{147}{ppm}$, en $2025 - 1750 = 275$ años. (c) La segunda
subida es mayor (147 frente a 90~ppm). Ritmos: $90 \div (2 \times 10^4) =
4.5 \times 10^{-3}$~ppm al año, frente a $147 \div 275 \approx 5.3 \times
10^{-1}$~ppm al año: unas 120 veces más rápida. (d) Que en 2019 era más
alto que en cualquier momento de al menos los dos últimos millones de años
(confianza alta).
\end{solution}

\begin{solution}{exo:g9:climate-past-and-present:12}
(a) Son casi iguales: 1.07 frente a \qty{1.06}{\degreeCelsius}. La
estimación central dice que la influencia humana explica aproximadamente
todo el calentamiento observado. (b) No: las causas naturales en conjunto
modificaron la temperatura \qty{0.1}{\degreeCelsius} como mucho, diez
veces menos que el calentamiento observado. (c) Al menos 66 posibilidades
sobre 100 de que el calentamiento causado por el ser humano esté entre
0.8 y \qty{1.3}{\degreeCelsius}.
\end{solution}

\begin{solution}{pb:g9:climate-past-and-present:1}
\textbf{1.} Porque es el periodo de referencia: las anomalías se miden a
partir de su propia media.

\textbf{2.} Más frío, en \qty{0.22}{\degreeCelsius}.

\textbf{3.} Más cálido, en \qty{0.98}{\degreeCelsius}.

\textbf{4.} $0.98 - (-0.22) = \qty{1.20}{\degreeCelsius}$.

\textbf{5.} $427 - 316 = \qty{111}{ppm}$.

\textbf{6.} $111 \div 316 \times 100 \approx 35~\%$.

\textbf{7.} $427 - 300 = \qty{127}{ppm}$.

\textbf{8.} El dióxido de carbono absorbe una parte de la luz infrarroja
que emite la superficie caliente y devuelve una parte hacia abajo. Con más
dióxido de carbono en el aire, vuelve más energía a la superficie, que se
calienta.

\textbf{9.} $2025 - 1993 = 32$ años; $32 \times 3.2 \approx
\qty{102}{mm}$, unos \qty{10}{cm}.

\textbf{10.} El agua del \omterm{def:g2:water-on-earth:ocean}{mar}, al calentarse, se dilata; el hielo
continental (\omterm{def:g8:ice-on-earth:glacier}{glaciares}, \omterm{def:g8:ice-on-earth:ice-sheet}{inlandsis}) se funde y va a parar al \omterm{def:g2:water-on-earth:ocean}{mar}.

\textbf{11.} Sí, del mismo orden: unos \qty{1.2}{\degreeCelsius} frente a
\qty{1.09}{\degreeCelsius}. Difieren un poco porque los periodos no son
los mismos (1880--1899 en lugar de 1850--1900; 2015--2024 en lugar de
2011--2020: los años más recientes fueron más cálidos), y porque el IPCC
combina varios registros, no solo el de la NASA.

\textbf{12.} ``Desde finales del siglo XIX, la superficie de la Tierra se
ha calentado unos \textbf{\qty{1.2}{\degreeCelsius}}''.
\end{solution}
